\documentclass[12pt]{article}
\usepackage[utf8]{inputenc}
\usepackage[spanish]{babel}
\usepackage{amsmath}
\usepackage{amssymb}

\title{Homología de grupos lineales sobre anillos de Dedekind}
\author{Abraham Rojas Vega}
\date{VII SEMINARIO INTERNACIONAL DE MATEMÁTICAS PERÚ-BRASIL 2025}

\begin{document}

\maketitle

\begin{abstract}
La homología de grupos lineales sobre anillos conmutativos es un campo que entrelaza la teoría de grupos, la topología algebraica y la teoría de números. Esta charla presenta una metodología para el cálculo de la segunda homología del grupo lineal especial sobre un anillo de Dedekind $A$, denotada como $H_2 (\operatorname{SL}_2 (A))$.

Nuestra aproximación se fundamenta en la teoría de árboles de Serre, la cual permite construir un complejo topológico sobre el que actúa el grupo $\operatorname{SL}_2 (A)$. Utilizamos la secuencia de Mayer-Vietoris, una herramienta de la topología algebraica, para descomponer el problema del cálculo de $H_2 (\operatorname{SL}_2 (A))$ en la homología de ciertos grupos de congruencia, que son más sencillos de analizar.

Esta charla estará enfocada en los aspectos conceptuales de la metodología, evitando detalles muy técnicos a fin de facilitar su comprensión para los alumnos. Se trata de un trabajo conjunto con Bruno Reis Ramos y Thiago Verissimo \cite{principal}.

\end{abstract}

\begin{thebibliography}{9}
\small
\bibitem{reference1}
Mirzaii, B., Ramos, B. R., Verissimo, T.
\textit{The second integral homology of $\operatorname{SL}_2(\mathbb{Z}[1/n])$}.
arXiv:2503.12190, 2025.

% \bibitem{reference1}
% Mirzaii, B., Ramos, B. R., Verissimo, T.
% \textit{Abelianization of $\operatorname{SL}_2$ over Dedekind domains of arithmetic type}.
% arXiv:2506.12638, 2025.

\bibitem{principal}
Ramos, B. R., Rojas Vega, A., Verissimo, T.
\textit{The Second Homology of $\operatorname{SL}_2$ over Dedekind Domains of Arithmetic Type}.
En preparación, 2025.

\bibitem{trees}
Serre, J.-P.
\textit{Trees}.
Springer, Berlin, 1980.

\end{thebibliography}
\end{document}

